\chapter*{Concluzii}
\addcontentsline{toc}{chapter}{Concluzii}

Lucrarea a prezentat un sistem complet de predicție a probabilității de câștig 
în meciuri de fotbal, combinând un model LightGBM calibrat cu o aplicație web 
full-stack funcțională, integrată cu datele Ligii 1 România prin API-Football.

Modelul antrenat pe \textbf{12.2 milioane de evenimente} din \textbf{3.464 de 
meciuri} StatsBomb obține un \textbf{RPS de 0.1431}, un \textbf{Brier Score de 
0.1472} și o \textbf{eroare de calibrare de 0.0309}, sub pragul academic de 
0.05, ceea ce îl face potrivit pentru uz tactic real \cite{Robberechts2019}. 
Concret, dacă modelul estimează $P = 70\%$ pentru victoria gazdei, echipa 
respectivă câștigă în $70\% \pm 3.09\%$ din situații similare.

Selecția LightGBM față de Gradient Boosting reflectă nu doar performanța 
competitivă (diferență sub 1\% pe toate metricile), ci și eficiența 
computațională net superioară: o serie completă de antrenare la 500 de meciuri 
s-a încheiat în 10 minute față de peste 17 ore pentru Gradient Boosting, un 
raport de aproximativ 100:1 care face LightGBM singurul model practic pentru 
reantrenare periodică pe volume crescânde de date.

Experimentele comparative confirmă două concluzii metodologice clare. 
Caracteristicile derivate de tip momentum, context temporal și context de joc aduc 
îmbunătățiri consistente pentru modelele de ansamblu, în timp ce Regresia 
Logistică nu poate exploata relațiile non-liniare introduse de aceste 
caracteristici, comportament așteptat, dat fiind că ipoteza sa de liniaritate 
este incompatibilă cu dinamica temporală a unui meci. Totodată, performanța se 
stabilizează în intervalul 200--500 de meciuri, sugerând că modelele de 
ansamblu converg rapid spre un nivel stabil fără a necesita volume foarte mari 
de date.

Aplicația dezvoltată pune aceste capabilități la dispoziția oricărui utilizator 
printr-o interfață web accesibilă: clasamentul și meciurile Ligii 1 România 
(sezon 2023) sunt preluate automat din API-Football, iar utilizatorul poate 
genera predicții probabilistice pentru orice meci finalizat printr-un singur 
click, predicțiile fiind salvate în istoricul său personal.

\textbf{Limitări.} Modelul operează pe baza unui snapshot al stării jocului 
la un moment dat, nu pe un flux continuu de evenimente. Nu sunt luați în 
considerare factori contextuali precum accidentările, condițiile meteo sau 
suspendările. Datele de antrenare acoperă exclusiv competițiile majore europene 
din StatsBomb, ceea ce poate afecta performanța pe ligi mai puțin reprezentate, 
inclusiv Liga 1 România, pe care aplicația o folosește în producție. 
Totodată, planul gratuit al API-Football limitează accesul la sezonul 2023, 
restricționând aplicabilitatea la date istorice.

\textbf{Direcții viitoare.} Arhitectura sistemului este proiectată explicit 
pentru a facilita extinderea naturală către predicție în timp real: înlocuirea 
sursei de date istorice cu un flux live de evenimente nu necesită modificări 
structurale majore în backend. Alte extensii posibile includ antrenarea 
incrementală pe date din Liga 1 pentru îmbunătățirea performanței locale, 
predicție minut-cu-minut pe toată durata meciului, incorporarea metricii 
Expected Threat (xT) ca o caracteristică suplimentară \cite{Singh2019} și explorarea 
modelelor de tip LSTM pentru capturarea dependențelor temporale dintre 
evenimentele unui meci.